% ===========================================================================
%  5 — Arquitetura de software
% ===========================================================================
\chapter{Arquitetura de \emph{software}}
\label{cap:arqsw}

Este capítulo descreve a organização interna do \emph{software}: que peças
existem, o que cada uma faz, como conversam e por que a fronteira está onde
está. A arquitetura de \emph{nuvem} --- onde essas peças rodam --- é o
Capítulo~\ref{cap:arqnuvem}.

\section{Visões arquiteturais}

O sistema é descrito por cinco visões complementares, no espírito do modelo
4+1~\cite{kruchten1995views}. Nenhuma delas é suficiente isoladamente.

\begin{tabularx}{\linewidth}{@{}L{30mm}L{40mm}X@{}}
\toprule
\headrow \thd{Visão} & \thd{Notação} & \thd{Onde está}\\
\midrule
Lógica & C4 níveis 1--3; UML: classes, objetos e o perfil da plataforma &
  Seções~\ref{sec:c4l1} a~\ref{sec:c4l4}; \ref{sec:uml-perfil} e~\ref{sec:uml-logica}\\
\zebra Processos & UML: sequência, estados, atividades, comunicação, tempo e visão
  geral de interação & Seções~\ref{sec:fluxos} e~\ref{sec:uml-processos};
  Cap.~\ref{cap:contratos}\\
Desenvolvimento & Estrutura do \emph{monorepo}; UML: pacotes, componentes e
  estrutura composta & Seções~\ref{sec:monorepo} e~\ref{sec:uml-desenvolvimento}\\
\zebra Física / implantação & Diagramas de implantação e de rede; UML: implantação
  com artefatos & Cap.~\ref{cap:arqnuvem}; Seção~\ref{sec:uml-fisica}\\
Cenários & Casos de uso & Cap.~\ref{cap:casosdeuso}; Seção~\ref{sec:uml-cenarios}\\
\bottomrule
\end{tabularx}

O C4 é a notação da arquitetura e continua sendo. A Seção~\ref{sec:uml} acrescenta
os catorze tipos de diagrama da UML 2.5.1, cada um onde diz algo que o C4 não tem
notação para dizer --- tipo, estado, prazo, instância, interface, artefato
(\adrr{ADR-122}).

\section{Estilo arquitetural e princípios}

A plataforma é, deliberadamente, um \textbf{monólito modular} de \emph{backend}
(uma API NestJS com módulos de domínio bem separados) servindo uma
\textbf{constelação de artefatos estáticos} (as experiências). Não é
microsserviços, e não é por acaso.

\begin{decisao}
\textbf{Um serviço de \emph{backend}, muitos artefatos de \emph{frontend}}
(\adrr{ADR-01}, \adrr{ADR-02}). Microsserviços resolvem escala organizacional
--- times independentes que precisam implantar sem coordenação. Aqui há uma
pessoa. O que precisa escalar independentemente é a \emph{entrega} das
experiências, e isso foi resolvido tirando-as do plano de execução:
\emph{bundles} estáticos numa CDN escalam sem que exista um serviço.
A modularidade que os microsserviços dariam por processo é obtida por
\emph{módulo NestJS} e por convenção de prefixo de tabela --- com a vantagem de
um \cd{JOIN} continuar sendo um \cd{JOIN}, e não uma chamada de rede.
\end{decisao}

Os princípios que atravessam todo o código:

\begin{description}
  \item[Validação na fronteira, sempre.] Todo \emph{input} externo passa por
    Zod~\cite{zod2026}: corpo, \emph{query}, \emph{postMessage} recebido e as
    variáveis de ambiente --- estas no \emph{boot}, de modo que configuração
    errada derruba o processo na subida, não no meio de uma requisição de estudante.
  \item[Uma porta por recurso crítico.] O banco só é acessado por
    \cd{PrismaService.comTenancy}; a IA só por \arq{src/ai}; o
    \emph{design system} só pela ponte \arq{packages/ui}; a entrada de pessoas e
    turmas só por \arq{src/rostering} --- que é também o único lugar da API que
    conhece o OneRoster, do mesmo jeito que só o \emph{provider} conhece o
    fornecedor de IA (\adrr{ADR-28}). Recurso com duas portas tem duas chances de
    esquecer a tranca.
  \item[Regra de domínio pura, casca fina.] Cada módulo de experiência separa
    \arq{catalogo.ts} (dados), \arq{regras.ts} (funções puras, testáveis sem
    banco) e um \emph{service} que persiste. O \emph{controller} é só a casca
    HTTP.
  \item[Falhar alto, ou falhar fechado.] Incompatibilidade de protocolo falha
    alto e visível; rota cuja proteção depende de configuração ausente responde
    401, nunca fica aberta.
  \item[O tipo carrega a regra.] Quando uma regra pode ser expressa no sistema de
    tipos, ela é --- porque comentário não é controle
    (Seção~\ref{sec:c4l4}).
\end{description}

\section{C4 nível 1 --- Contexto do sistema}
\label{sec:c4l1}

\begin{figure}[H]
\centering
\fitwidth{%
\begin{tikzpicture}[x=1cm, y=1cm]

  % --- pessoas
  \node[c4person] (estudante) at (1.9,0)  {\textbf{Estudante}\desc{1º EF ao 3º EM}};
  \node[c4person] (prof)  at (6.0,0)  {\textbf{Professor}\desc{Conduz a aula}};
  \node[c4person] (gest)  at (10.1,0) {\textbf{Gestor escolar}\desc{Mantém usuários}};
  \node[c4person] (admin) at (14.2,0) {\textbf{Admin do contrato}\desc{Rede}};

  % --- sistema central
  \node[c4system, text width=76mm, minimum height=26mm] (sis) at (8.0,-4.4)
    {\textbf{Plataforma Unooria}\desc{Hub multi-tenant que autentica por papel,
     embeda experiências imersivas versionadas, corrige e pontua no servidor,
     coleta telemetria de aprendizagem e produz leitura de turma --- com
     isolamento lógico por escola}};

  % --- sistemas externos
  \node[c4systemext, text width=26mm, minimum height=18mm] (idp) at (1.6,-4.4)
    {\textbf{WSO2 Identity Server}\desc{Identidade central, vários contratos}};
  \node[c4systemext, text width=26mm, minimum height=18mm] (vertex) at (14.4,-4.4)
    {\textbf{Vertex AI --- Gemini}\desc{Texto e imagem generativos}};
  \node[c4systemext, text width=32mm, minimum height=14mm] (ci) at (3.4,-8.6)
    {\textbf{GitHub Actions}\desc{Publica \emph{bundles} e manifestos}};
  \node[c4systemext, text width=32mm, minimum height=14mm] (tf) at (12.6,-8.6)
    {\textbf{Terraform (repo separado)}\desc{Provisiona a \emph{stack} do contrato}};
  % --- o vendor da Ajuda com avatar (ADR-132): fora do Google, cobra por minuto
  \node[c4systemext, text width=30mm, minimum height=14mm] (avatar) at (8.0,-8.6)
    {\textbf{LiveAvatar (HeyGen)}\desc{Escuta, voz e vídeo do atendente da Ajuda}};

  % --- relações
  \draw[c4rel] (estudante.south) -- (4.4,-3.1)
    node[cxlabel, pos=0.42, left=1pt]{realiza missões};
  \draw[c4rel] (prof.south) -- (6.6,-3.1)
    node[cxlabel, pos=0.40, right=1pt]{prepara, conduz, lê};
  \draw[c4rel] (gest.south) -- (9.6,-3.1)
    node[cxlabel, pos=0.40, left=1pt]{administra};
  \draw[c4rel] (admin.south) -- (11.6,-3.1)
    node[cxlabel, pos=0.42, right=1pt]{lê o agregado};
  \draw[c4reldash] (idp.east) -- (sis.west)
    node[cxlabel, midway, above=3pt]{OIDC};
  \draw[c4relbi] (sis.east) -- (vertex.west)
    node[cxlabel, midway, above=3pt]{HTTPS + ADC};
  \draw[c4rel] (ci.north) -- (5.6,-5.9)
    node[cxlabel, pos=0.5, left=1pt]{publica};
  \draw[c4reldash] (tf.north) -- (10.6,-5.9)
    node[cxlabel, pos=0.5, right=1pt]{provisiona};
  \draw[c4relbi] (sis.south) -- (avatar.north)
    node[cxlabel, midway, right=1pt]{sala; \texttt{/ajuda/llm}};

\end{tikzpicture}}
\caption[C4 nível 1 --- contexto]{C4 nível 1: contexto. Quatro papéis humanos
  mutuamente exclusivos, um provedor de identidade compartilhado entre contratos,
  dois serviços externos de IA, e dois atores de automação. O segundo serviço de
  IA, desde 2026-10-02, é o LiveAvatar da Ajuda com avatar (\adrr{ADR-132}): a
  relação é de mão dupla porque a API pede a sala de vídeo ao vendor e o vendor
  chama de volta o LLM da API (\rota{/ajuda/llm}); e é o único externo que o
  navegador também alcança direto, pela sala de vídeo (WebRTC). O
  Terraform aparece porque a fronteira entre os dois repositórios é um contrato
  explícito, não um detalhe de operação.}
\label{fig:c4l1}
\end{figure}

\begin{nota}
A dependência externa \emph{não substituível em tempo de execução} é o Vertex AI,
e mesmo assim sua indisponibilidade não derruba a plataforma: as experiências que
não usam IA (exp04ef1b, app-70) seguem íntegras, o \emph{feed} da exp08ef1b continua
sendo curado a partir de conteúdo já gerado, e as ferramentas de professor
degradam para texto fixo (app-91) ou recusam a geração com mensagem clara
(app-90). Essa contenção de acoplamento é deliberada (\rf{RNF-12}).
O LiveAvatar é ainda mais contido: fora do ar, só a Ajuda falha
(\rota{POST /ajuda/sessoes} responde 502 \cd{AVATAR\_FALHOU}), e sem as variáveis
dele a API sobe normal e a Ajuda responde 503 (\adrr{ADR-132}).
\end{nota}

\section{C4 nível 2 --- Contêineres}
\label{sec:c4l2}

O nível de contêineres é onde a assimetria da arquitetura aparece. Há
\emph{cinco} contêineres de execução e um número crescente de \emph{artefatos}
que não são contêineres --- e tratá-los como se fossem seria o erro que este
desenho existe para evitar. O quinto, desde 2026-09-30, é o agente de IA do
professor no Vertex AI Agent Engine (\adrr{ADR-126}): o único contêiner que
\textbf{só a API chama}, e o único escrito em Python.

\begin{figure}[H]
\centering
\fitpage{%
\begin{tikzpicture}[x=1cm, y=1cm]

  % --- fronteira do sistema
  \draw[c4zone] (0.6,-1.6) rectangle (15.4,-17.4);
  \node[c4boundarylbl, fill=unooSurface] at (1.0,-1.45)
    {Plataforma Unooria --- \emph{stack} de um contrato};

  % --- pessoas
  \node[c4person, text width=24mm, minimum height=12mm] (estudante) at (3.0,-0.2)
    {\textbf{Estudante}};
  \node[c4person, text width=24mm, minimum height=12mm] (prof) at (7.0,-0.2)
    {\textbf{Professor}};
  \node[c4person, text width=26mm, minimum height=12mm] (adm) at (11.4,-0.2)
    {\textbf{Gestor / Admin}};

  % --- navegador
  \draw[iframe] (8.9,-2.4) rectangle (14.9,-6.1);
  \node[iframelbl] at (9.1,-2.3) {\texttt{iframe} --- \emph{sandbox}};

  \node[c4container, text width=34mm, minimum height=17mm] (hub) at (4.3,-4.0)
    {\textbf{Hub}\tech{Next.js App Router}\desc{Sessão, roteamento por papel,
     \emph{chrome} institucional, montagem do \emph{embed}, BFF, JWKS e
     assinatura de \emph{token}}};

  \node[bundle, text width=34mm, minimum height=17mm] (exp) at (11.9,-4.2)
    {\textbf{Experiência}\tech{Vite + React, estático}\desc{Uma por missão.
     Roda no navegador do estudante; nunca é um processo do servidor}};

  \node[c4component, text width=24mm, minimum height=11mm] (sdk) at (11.9,-7.3)
    {\textbf{\texttt{@edu/sdk}}\tech{TypeScript}\desc{Ponte, telemetria, API,
     progresso}};

  \node[c4container, text width=36mm, minimum height=19mm] (api) at (4.3,-8.6)
    {\textbf{API}\tech{NestJS + Prisma}\desc{Autorização por papel,
     \emph{tenancy}, regra de domínio por experiência, gabarito, correção,
     \emph{gateway} de IA, ingestão de telemetria}};

  \node[c4db, text width=27mm, minimum height=19mm] (pg) at (8.9,-8.6)
    {\textbf{PostgreSQL}\tech{Cloud SQL, HA + PITR}\desc{Único banco, único
     \emph{schema}; RLS por escola}};

  \node[c4queue, text width=26mm, minimum height=13mm] (ps) at (13.2,-9.1)
    {\textbf{Pub/Sub}\desc{Canal analítico de telemetria}};

  \node[c4artifact, text width=28mm, minimum height=12mm] (buckets) at (2.6,-12.2)
    {\textbf{Object Storage}\desc{\emph{bundles}, \emph{assets} públicos,
     planos privados}};

  \node[c4systemext, text width=26mm, minimum height=12mm] (vertex) at (7.4,-12.2)
    {\textbf{Vertex AI}};
  \node[c4containerext, text width=26mm, minimum height=12mm] (idp) at (11.6,-12.2)
    {\textbf{WSO2 IS}\desc{externo à \emph{stack}}};

  % --- o agente de IA do professor (ADR-126/127): so a API o chama
  \node[c4container, text width=30mm, minimum height=15mm] (engine) at (7.6,-15.4)
    {\textbf{Agente do professor}\tech{Python + ADK, Agent Engine}\desc{Sessão,
     ferramentas, guarda de escopo e título}};
  \node[c4db, text width=30mm, minimum height=15mm] (ds) at (12.4,-15.6)
    {\textbf{Vertex AI Search}\tech{dois \emph{data stores}, \texttt{global}}\desc{habilidades
     (estruturado) e acervo (PDFs)}};

  % --- o vendor da Ajuda com avatar (ADR-132): fora da stack e fora do Google
  \node[c4systemext, text width=26mm, minimum height=12mm] (avatar) at (2.6,-16.4)
    {\textbf{LiveAvatar}\desc{HeyGen; vídeo da Ajuda}};

  % --- relações
  \draw[c4rel] (estudante.south) -- (3.6,-3.1) node[cxlabel, pos=0.45, left=1pt]{HTTPS};
  \draw[c4rel] (prof.south) -- (5.0,-3.1) node[cxlabel, pos=0.45, right=1pt]{HTTPS};
  \draw[c4rel] (adm.south) -- (5.4,-3.3);
  \draw[c4relbi] (hub.east) -- (exp.west)
    node[cxlabel, pos=0.5, above]{\texttt{postMessage} v1}
    node[cxlabel, pos=0.5, below]{\emph{targetOrigin} explícito};
  \draw[c4rel] (exp.south) -- (sdk.north);
  \draw[c4rel] (sdk.south west) -- (api.north east)
    node[cxlabel, pos=0.42, above right=-2pt]{HTTPS + \emph{Bearer}};
  \draw[c4rel] (hub.south) -- (api.north) node[cxlabel, midway, right=2pt]{BFF};
  \draw[c4relbi] (api.east) -- (pg.west) node[cxlabel, midway, above]{SQL + RLS};
  \draw[c4rel] (api.south east) -- (ps.south west)
    node[cxlabel, pos=0.72, below]{publica};
  \draw[c4rel] (ps.north west) -- (5.9,-7.4)
    node[cxlabel, pos=0.62, above]{\emph{push} \texttt{/telemetry/ingest}};
  \draw[c4rel] (api.south) -- (buckets.north east)
    node[cxlabel, pos=0.5, left=1pt]{assets, planos};
  \draw[c4rel] (api.south) -- (vertex.north);
  \draw[c4rel] (5.4,-10.5) -- (5.4,-15.4) -- (engine.west)
    node[cxlabel, pos=0.3, left=1pt]{\texttt{:streamQuery}, ADC};
  \draw[c4rel] (engine.north) -- (vertex.south)
    node[cxlabel, midway, right=1pt]{Gemini};
  \draw[c4rel] (engine.east) -- (engine.east -| ds.west) node[cxlabel, midway, above]{busca};
  \draw[c4relbi] (4.8,0 |- api.south) -- (4.8,-14.0) -- (2.6,-14.0)
    node[cxlabel, pos=0.6, above]{sala; \texttt{/ajuda/llm}} -- (avatar.north);
  \draw[c4reldash] (hub.west) -- (0.9,-4.0) -- (0.9,-12.2) -- (idp.west)
    node[cxlabel, pos=0.86, above]{OIDC};
  \draw[c4reldash] (buckets.north) -- (2.6,-6.6) -- (9.2,-6.6)
    node[cxlabel, pos=0.72, above]{CDN serve o \emph{bundle}};

\end{tikzpicture}}
\caption[C4 nível 2 --- contêineres]{C4 nível 2: contêineres. Só o hub, a API, o
  banco, o Pub/Sub e o agente de IA do professor são peças em execução; o agente
  roda no Vertex AI Agent Engine e só a API o chama --- o navegador nunca fala
  com ele nem com os \emph{data stores} (\adrr{ADR-126}). O LiveAvatar, da Ajuda
  com avatar, é externo à \emph{stack}: a API abre e fecha a sala e o vendor chama
  de volta \rota{/ajuda/llm} pela URL pública da API; o vídeo vai direto do
  vendor ao navegador (WebRTC), sem passar pela \emph{stack}; a chave da conta,
  o \emph{prompt} e a base nunca saem da API (\adrr{ADR-132}). A experiência é um \emph{artefato}
  servido pela CDN e executado no navegador; a SDK é biblioteca dentro dele. A
  linha tracejada laranja é a fronteira do \emph{sandbox} do \cd{iframe} --- a
  fronteira de segurança mais importante do desenho de \emph{frontend}.}
\label{fig:c4l2}
\end{figure}

\begin{legendadiag}
\amostra{c4container} contêiner em execução \quad
\amostra{bundle} artefato estático \quad
\amostra{c4component} biblioteca \quad
\amostra{c4db} banco \quad
\amostra{c4queue} fila \quad
\amostra{c4artifact} armazenamento de objetos \quad
\amostra{c4systemext} externo
\end{legendadiag}

\subsection{Catálogo de contêineres e artefatos}

\begin{tabularx}{\linewidth}{@{}L{26mm}L{20mm}L{22mm}X@{}}
\toprule
\headrow \thd{Peça} & \thd{Tipo} & \thd{Tecnologia} & \thd{Responsabilidade e limite}\\
\midrule
\endhead
\arq{apps/hub} & Contêiner & Next.js 16, App Router &
  Autenticar, roteirizar por papel, desenhar \emph{todo} o \emph{chrome}
  institucional, montar o \emph{embed}, assinar \emph{token} de experiência e
  publicar JWKS. \textbf{Não} contém regra de domínio de experiência.\\
\zebra \arq{apps/api} & Contêiner & NestJS, Prisma &
  Autorizar, garantir \emph{tenancy}, aplicar regra de domínio, guardar
  gabarito, corrigir, intermediar IA, ingerir telemetria. É o único que fala com
  o banco.\\
\arq{apps/experiencias/*} & Artefato & Vite, React &
  Interface e interação de uma missão. \textbf{Nunca} guarda gabarito, nem
  calcula nota que valha, nem persiste \emph{token}.\\
\zebra \arq{packages/sdk} & Biblioteca & TypeScript puro, Zod &
  Contrato \cd{postMessage}, cliente HTTP autenticado, telemetria em lote,
  progresso. É a \emph{única} forma legítima de uma experiência falar com a
  plataforma.\\
\arq{packages/ui} & Biblioteca & TypeScript, React &
  Ponte para o \emph{design system}: marca \cd{'use client'}, desfaz colisão de
  nomes entre seções e preserva \emph{tree-shaking}.\\
\zebra PostgreSQL & Contêiner gerenciado & Cloud SQL &
  Único banco, único \emph{schema} lógico. Guarda tudo, com RLS forçada.\\
Pub/Sub & Contêiner gerenciado & Pub/Sub &
  Amortecedor \emph{durável} do pico de telemetria. Entrega \emph{at least
  once} por \emph{push} autenticado.\\
\zebra Object Storage & Artefato & Cloud Storage &
  Três \emph{buckets} com regimes diferentes: \emph{bundles} (público, imutável),
  \emph{assets} (público, CDN), planos (\textbf{privado}, versionado).\\
\arq{apps/\B agente-professor} & Contêiner gerenciado & Python 3.12, Google ADK, Vertex AI Agent Engine &
  O agente de IA do professor: instrução, ferramentas, guarda de escopo, título
  e a sessão (histórico). \textbf{Só a API o chama}; não conhece papel, escola
  nem quota --- a posse da conversa é do Postgres (\adrr{ADR-126}).\\
\zebra \emph{Data stores} & Contêiner gerenciado & Vertex AI Search, \cd{global} &
  \cd{agente-professor-\B habilidades} (estruturado, uma linha por habilidade) e
  \cd{agente-professor-\B acervo} (PDFs, \emph{chunking} por \emph{layout}), alimentados
  do \emph{bucket} \cd{<projeto>-\B agente-professor} por \arq{ingestao/}
  (\adrr{ADR-127}).\\
\bottomrule
\end{tabularx}

\begin{regra}
\textbf{Um banco, um \emph{schema}.} Nem banco por experiência --- elas não são
isoladas (o perfil da Vila alimenta a largada do Instanooria, e o PostgreSQL não
faz \cd{JOIN} entre bancos; além disso, um \emph{pool} do Prisma por experiência
estouraria \cd{max\_connections}) --- nem \emph{schema} por experiência: a única
coisa que um \emph{schema} daria além do prefixo é fronteira de \cd{GRANT}, e
existe \emph{uma} \emph{role} de aplicação. A separação é por
\textbf{prefixo de tabela}: o código da experiência (\cd{exp08ef1b\_},
\cd{app\_70\_} --- \adrr{ADR-110}) para tabela de experiência, sem
prefixo para tabela de plataforma (Capítulo~\ref{cap:dados}).
\end{regra}

\section{C4 nível 3 --- Componentes}
\label{sec:c4l3}

\subsection{Componentes da API}

A API é organizada em módulos NestJS. A Figura~\ref{fig:c4l3api} separa-os em
quatro faixas, e a faixa a que um módulo pertence determina o que ele pode
conhecer. Desde 2026-09-02 a faixa transversal tem um quinto componente que
a figura não desenha, \cd{observabilidade/} (Capítulo~\ref{cap:observabilidade}):
o \emph{middleware} que abre o contexto da requisição antes dos \emph{guards}
(o \cd{TenancyInterceptor} passou a \emph{preencher} a escola nesse mesmo
objeto, em vez de abri-lo), o \emph{logger} pino com \emph{redaction}, o
\cd{FiltroDeExcecao} global, o \emph{sampler} por rota e o \arq{otel.ts} ---
que é a \textbf{primeira importação} do \arq{main.ts}, porque a
instrumentação automática do OpenTelemetry precisa alterar \cd{http},
\cd{express}, \cd{undici} e o Prisma antes de o Nest os carregar.

Há um segundo componente que a figura não desenha, de 2026-09-30:
\arq{agente-professor/}, o módulo do agente de IA do professor
(\adrr{ADR-126}). Ele é de \textbf{plataforma}, não de experiência --- o agente é
um produto sem experiência, trancado por \cd{@Produto(\B"agente-\B professor")} ---, e
usa três portas: \cd{comTenancy} para a posse da conversa
(\cd{agente\_\B professor\_\B conversas}), o \cd{QuotaService} de \arq{ai/} para o
saldo e o consumo, e a \cd{AgenteProfessorPort}, a segunda porta de \arq{ai/},
que só \arq{ai/\B agent-engine.\B client.ts} implementa. De \arq{curriculo/} ele lê
\cd{habilidade\B Por\B Codigo}, como \arq{evidencias/}, para conferir os códigos que o
agente cita (\adrr{ADR-127}). O par Mermaid do
Apêndice~\ref{ap:mermaid} traz as relações dele.

O terceiro, de 2026-10-02, é \arq{ajuda/}, a Ajuda com avatar
(\adrr{ADR-132}): o atendente em vídeo do menu da conta, para todo papel menos o
estudante. Também é de \textbf{plataforma} e produto sem experiência
(\cd{@Produto(\B"ajuda")}), e tem \textbf{dois} \emph{controllers} com portas de
entrada diferentes. \arq{ajuda.\B controller.ts} é o lado do hub (abrir e encerrar
a sessão), atrás de \cd{OidcGuard}, \cd{RolesGuard} e \cd{DisponibilidadeGuard}
--- o encerrar escapa do interruptor (\cd{@SemDisponibilidade}), porque desligar
o produto no meio da conversa não pode deixar a sala cobrando.
\arq{ajuda-llm.\B controller.ts} é o lado do \emph{vendor}: o \emph{endpoint}
OpenAI-compatível que o LiveAvatar chama a cada fala, \textbf{fora} dos
\emph{guards} de pessoa, trancado por duas chaves (o \emph{Bearer} do
\emph{vendor} e o \emph{token} curto da sessão, que viaja dentro da mensagem
\cd{system}). Ele usa duas portas de \arq{ai/}: o \cd{AiProvider}, pela
\cd{conversarEmFluxo} --- a única geração em \emph{stream} do \emph{provider}
---, e a \cd{AvatarAoVivoPort} (\emph{token} \cd{AVATAR\_AO\_VIVO}), a terceira
porta de \arq{ai/}, que só \arq{ai/\B liveavatar.\B client.ts} implementa e que só
abre e fecha a sala de vídeo; é porta própria, e não método do \cd{AiProvider},
porque é reserva de um recurso cobrado por minuto, com ciclo de vida. Não passa
pelo \cd{QuotaService}: o custo é limitado em minutos (\cd{ajuda\_\B limites},
editada pelo \cd{admin} em \arq{recursos/\B limites-da-ajuda.\B controller.ts}) e
medido em \cd{ajuda\_\B sessoes}, as duas fora de \cd{comTenancy} e consultadas
pelo \cd{usuario\_id} do \emph{token} --- o \cd{gestor\_contrato} e o \cd{admin}
não têm escola. A base de conhecimento é gerada
(\arq{ajuda/\B base/\B base-da-ajuda.\B gerada.ts}) e nunca editada à mão. As
relações dele também estão só no par Mermaid.

O quarto, de 2026-10-06, é \arq{podcast/}, o podcast de recapitulação
(\adrr{ADR-150}). Também é de \textbf{plataforma}: as rotas
(\rota{/podcast/:expId/*}, do professor e da gestão), o motor do áudio
(\cd{generateSpeech} por trecho e a emenda do MP3 em \arq{podcast/\B mp3.ts}), a
escrita das falas do tema pela IA (com o \cd{QuotaService}) e a tabela
\cd{podcast\_\B roteiros}. Não conhece missão nenhuma pelo nome: cada missão com
podcast --- \arq{a-selecao/}, \arq{jornal-escola/}, \arq{quem-mandou-isso/} e,
desde a \adrr{ADR-151}, \arq{juri-estadual/}; desde a \adrr{ADR-156},
\arq{atividades/} (o Unoovista, que grava as falas no Ativar) e \arq{unoogpt/}
(a Missão Tempestade, o primeiro podcast fixo, sem vaga) ---
registra uma \cd{FonteDoPodcast} na subida (\cd{PodcastService.\B registrarFonte}),
no molde do \cd{registrarContexto} de \arq{evidencias/}, e a homologação dela
grava as falas aprovadas com \cd{gravarFalasDoPodcast}, função livre sobre o
\cd{tx}. O módulo é \cd{@Global()} e o último do \cd{app.module}, para que uma
missão nova com podcast não reordene o OpenAPI. As relações dele também estão
só no par Mermaid.

\begin{figure}[H]
\centering
\fitpage{%
\begin{tikzpicture}[x=1cm, y=1cm]
  % Grade (2026-09-25): cinco colunas de 32mm de passo (x = 1,85 / 5,05 / 8,25 /
  % 11,45 / 14,65) com caixas de 25,5mm e 3mm de vão entre elas; a faixa
  % transversal usa quatro colunas mais largas porque `TenancyInterceptor` não
  % cabe em 25,5mm. A versão anterior encavalava caixas (`salas/` sobre
  % `health/`, `recursos/` sobre `analise/`) e o rótulo do domínio sobre as
  % duas: typecheck nenhum pega isso --- só a página renderizada.
  % A infraestrutura fica ENTRE a plataforma e o domínio para que as duas
  % faixas que a usam cheguem a ela sem atravessar a outra; a coluna 2 da
  % infraestrutura fica VAZIA de propósito: é por ela que os painéis sobem até
  % `equipes/`. Os rótulos das faixas vão à direita porque a esquerda é a
  % pista das setas.

  \draw[c4zone] (0,0) rectangle (16.5,-15.95);
  \node[c4boundarylbl, fill=unooSurface] at (0.4,0.15) {Contêiner: API (NestJS)};

  % ---- faixa 1: transversal
  \node[camadalbl, anchor=east] at (16.1,-0.5) {Transversal --- toda rota passa por aqui};
  \node[c4component, text width=31mm, minimum height=16mm] (guard) at (2.125,-1.55)
    {\texttt{OidcGuard}\desc{valida JWT por JWKS}};
  \node[c4component, text width=31mm, minimum height=16mm] (roles) at (6.208,-1.55)
    {\texttt{RolesGuard}\desc{restringe por papel}};
  \node[c4component, text width=31mm, minimum height=16mm] (tenint) at (10.292,-1.55)
    {\texttt{TenancyInterceptor}\desc{escola no \texttt{AsyncLocalStorage}}};
  \node[c4component, text width=31mm, minimum height=16mm] (env) at (14.375,-1.55)
    {\texttt{validarEnv}\desc{Zod no \emph{boot}}};

  % ---- faixa 2: plataforma
  \node[camadalbl, anchor=east] at (16.1,-2.95) {Plataforma --- serve a todas as experiências};
  \node[c4component, text width=25.5mm, minimum height=16mm] (auth) at (1.85,-4.0)
    {\texttt{auth/}\desc{\texttt{TokenService}}};
  \node[c4component, text width=25.5mm, minimum height=16mm] (usu) at (5.05,-4.0)
    {\texttt{usuarios/}\desc{listar, editar, senha}};
  \node[c4component, text width=25.5mm, minimum height=16mm] (expmod) at (8.25,-4.0)
    {\texttt{experiencias/}\desc{manifesto}};
  \node[c4component, text width=25.5mm, minimum height=16mm] (tel) at (11.45,-4.0)
    {\texttt{telemetry/}\desc{\emph{batch}, \emph{ingest}}};
  \node[c4component, text width=25.5mm, minimum height=16mm] (prog) at (14.65,-4.0)
    {\texttt{progress/}\desc{progresso genérico}};
  \node[c4component, text width=25.5mm, minimum height=16mm] (health) at (1.85,-5.95)
    {\texttt{health/}\desc{\texttt{healthz}, \texttt{readyz}}};
  \node[c4component, text width=25.5mm, minimum height=16mm] (analise) at (5.05,-5.95)
    {\texttt{analise/}\desc{turma e rede, CSV}};
  \node[c4component, text width=25.5mm, minimum height=16mm] (comum) at (8.25,-5.95)
    {\texttt{comum/}\desc{séries, escopo de turma}};
  \node[c4component, text width=25.5mm, minimum height=16mm] (rost) at (11.45,-5.95)
    {\texttt{rostering/}\desc{pessoas, \mbox{escolas}, turmas}};
  \node[c4component, text width=25.5mm, minimum height=16mm] (port) at (14.65,-5.95)
    {\texttt{portabilidade/}\desc{perfil portátil}};
  \node[c4component, text width=25.5mm, minimum height=16mm] (salas) at (1.85,-8.1)
    {\texttt{salas/}\desc{a sala: grade e liberação}};
  \node[c4component, text width=25.5mm, minimum height=16mm] (equipes) at (5.05,-8.1)
    {\texttt{equipes/}\desc{quem divide um computador}};
  \node[c4component, text width=25.5mm, minimum height=16mm] (recs) at (8.25,-8.1)
    {\texttt{recursos/}\desc{o interruptor da rede e a versão}};
  \node[c4component, text width=25.5mm, minimum height=16mm] (prosp) at (11.45,-8.1)
    {\texttt{prospeccao/}\desc{o site público, sem \emph{tenant}}};
  \node[c4component, text width=25.5mm, minimum height=16mm] (evid) at (14.65,-8.1)
    {\texttt{evidencias/}\desc{habilidade da BNCC por atividade}};

  % ---- faixa 3: infraestrutura (coluna 2 vazia: a pista das equipes)
  \node[camadalbl, anchor=east] at (16.1,-9.25) {Infraestrutura --- as portas únicas};
  \node[c4component, text width=25.5mm, minimum height=16mm] (prisma) at (1.85,-10.4)
    {\texttt{PrismaService}\desc{\texttt{comTenancy}: transação + RLS}};
  \node[c4component, text width=25.5mm, minimum height=16mm] (ai) at (8.25,-10.4)
    {\texttt{ai/}\desc{\texttt{AiProvider}, \texttt{QuotaService}}};
  \node[c4component, text width=25.5mm, minimum height=16mm] (store) at (11.45,-10.4)
    {\texttt{object-storage/}\desc{local $|$ GCS}};
  \node[c4component, text width=25.5mm, minimum height=16mm] (assets) at (14.65,-10.4)
    {\texttt{assets/}\desc{imagem gerada}};

  % ---- faixa 4: dominio por experiencia
  % Linha de cima alinhada à grade; a de baixo DESLOCADA meia coluna, para a
  % seta de cada caixa subir pelo vão entre duas da linha de cima.
  \node[camadalbl, anchor=east] at (16.1,-11.82) {Domínio por experiência};
  \node[c4component, text width=25.5mm, minimum height=19.5mm, fill=unooGreen!22, draw=unooGreen] (motor) at (1.85,-13.075)
    {\texttt{motor/} e mais 4\desc{geracao, dossie, atividades, perfis --- exp08ef1b}};
  \node[c4component, text width=25.5mm, minimum height=19.5mm, fill=unooGreen!22, draw=unooGreen] (utube) at (5.05,-13.075)
    {\texttt{unootube/}\desc{exp04ef1b}};
  \node[c4component, text width=25.5mm, minimum height=19.5mm, fill=unooGreen!22, draw=unooGreen] (bem) at (8.25,-13.075)
    {\texttt{bemestar/}\desc{app-91}};
  % planos/ ocupa a quarta coluna da linha de cima desde que o o-eco-no-grupo/
  % saiu (ADR-113): é o módulo que usa as mesmas DUAS portas.
  \node[c4component, text width=25.5mm, minimum height=19.5mm, fill=unooGreen!22, draw=unooGreen] (planos) at (11.45,-13.075)
    {\texttt{planos/}\desc{app-90}};
  \node[c4component, text width=27mm, minimum height=35.5mm, fill=unooGreen!22, draw=unooGreen, font=\scriptsize] (outros) at (14.73,-13.875)
    {\texttt{rastro-boato/}\\\texttt{manifesto-ia/}\\\texttt{quem-mandou-isso/}\\\texttt{foto-conta/}\\\texttt{planejamento/}};
  \node[c4component, text width=25.5mm, minimum height=13mm, fill=unooGreen!22, draw=unooGreen] (vila) at (3.45,-15.0)
    {\texttt{vila/}\desc{app-70}};
  \node[c4component, text width=25.5mm, minimum height=13mm, fill=unooGreen!22, draw=unooGreen] (ugpt) at (6.65,-15.0)
    {\texttt{unoogpt/}\desc{exp02em1b}};

  % ---- relações: transversal
  \draw[c4rel] (guard.east) -- (roles.west);
  % a escola do token desce pela margem esquerda, fora de toda caixa
  \draw[c4reldash] (tenint.south) -- (10.292,-2.6) -- (0.2,-2.6) -- (0.2,-10.4) -- (prisma.west);
  \node[cxlabel] at (6.208,-2.6) {escola do \emph{token}};

  % ---- relações: plataforma
  \draw[c4rel] (2.6,-7.3) -- (2.6,-7.05) -- (4.6,-7.05) -- (4.6,-7.3);
  \node[cxlabel] at (3.6,-7.05) {contagem};
  \draw[c4rel] (salas.south) -- (prisma.north);
  \draw[c4rel] (3.75,-8.9) -- (3.75,-9.25) -- (2.9,-9.25) -- (2.9,-9.6);
  % evidencias/ -> equipes/: pelo vão entre a 2a e a 3a linha da plataforma,
  % acima de recursos/ e prospeccao/, e descendo no terço direito de equipes/
  % (o terço esquerdo é a chegada da contagem da sala).
  \draw[c4rel] (evid.north) -- (14.65,-7.05) -- (5.6,-7.05) -- (5.6,-7.3);
  \node[cxlabel] at (10.05,-7.05) {gancho \texttt{aoMudarEquipes} + \texttt{consulta.ts}};

  % ---- relações: domínio (as que sobem até equipes/ passam pela coluna vazia)
  \draw[c4rel] (motor.north) -- (1.85,-11.2);
  \draw[c4rel] (3.45,-14.35) -- (3.45,-11.75) -- (2.9,-11.75) -- (2.9,-11.2);
  \draw[c4rel] (4.4,-12.1) -- (4.4,-8.9);
  \draw[c4rel] (6.65,-14.35) -- (6.65,-11.75) -- (5.0,-11.75) -- (5.0,-8.9);
  \draw[c4rel] (6.65,-11.75) -- (7.4,-11.75) -- (7.4,-11.2);
  \node[cxlabel, anchor=west] at (5.15,-9.25) {\texttt{consulta.ts}};
  \draw[c4rel] (bem.north) -- (8.25,-11.2);
  % planos usa as DUAS portas: uma saída sobe ao bucket, a outra desvia para a IA
  \draw[c4rel] (planos.north) -- (11.45,-11.2);
  \draw[c4rel] (11.45,-11.75) -- (9.2,-11.75) -- (9.2,-11.2);
  \fill[unooInkSoft] (11.45,-11.75) circle (1.5pt)
                     (6.65,-11.75) circle (1.5pt);
  \node[cxlabel] at (10.65,-11.47) {IA + quota};
  \node[cxlabel, anchor=west] at (11.55,-11.47) {\emph{bucket}};

\end{tikzpicture}}
\caption[C4 nível 3 --- componentes da API]{C4 nível 3: componentes da API. A
  faixa transversal roda em toda rota; a de plataforma serve a todas as
  experiências; a de infraestrutura contém as portas únicas (banco, IA,
  armazenamento); a de domínio tem um módulo por experiência ---
  \texttt{catalogo.ts} + \texttt{regras.ts} + \emph{service} + \emph{controllers}
  por papel ---, e nenhum módulo dessa faixa conhece outro. \texttt{equipes/}
  (\adrr{ADR-111}) é lido por \texttt{salas/} (\adrr{ADR-69}), que mostra a
  contagem na sala, e pelos painéis de \texttt{unootube/} e \texttt{unoogpt/},
  pela porta \texttt{equipes/consulta.ts}; \texttt{evidencias/}
  (\adrr{ADR-112}) lê a mesma porta e se registra no gancho
  \texttt{aoMudarEquipes}, e as missões o chamam pela porta \texttt{marcar}
  --- hoje só \texttt{quem-mandou-isso/} (\adrr{ADR-118}), que também lê
  \texttt{equipes/consulta.ts} e usa IA com quota no preparo. A caixa da
  direita reúne, sem setas, os módulos da exp06ef2b, da exp08ef4b, da
  exp05ef1b, da exp01ef1b e da app-92, na ordem em que aparecem; as setas da
  exp05ef1b estão no par Mermaid, para não cruzarem a faixa de plataforma.
  Todo módulo com dado passa por
  \texttt{comTenancy}: as setas até o \texttt{PrismaService} são uma amostra, e a
  fonte Mermaid do Apêndice~\ref{ap:mermaid} lista as relações por módulo.}
\label{fig:c4l3api}
\end{figure}

\begin{regra}
\begin{nota}
\arq{unoogpt/} é o módulo de domínio mais fundo da API: 3.285 linhas de motor
probabilístico, agente e regras puras, com subpastas (\arq{modelo/},
\arq{regras/}, \arq{agente/}) em vez do par \arq{catalogo.ts} +
\arq{regras.ts}. A granularidade veio do projeto de origem e foi \textbf{mantida
de propósito}: cada arquivo tem teste próprio, e colar arquivos bem fatorados
produziria colisão de nomes e falha mais difícil de localizar. O padrão do
Capítulo~\ref{cap:experiencias} continua valendo para experiência nova; este é um
porte.
\end{nota}

\textbf{\arq{recursos/} é o terceiro eixo, e ele não é de \emph{tenancy}}
(\adrr{ADR-87}). O catálogo (\arq{experiencias/}) diz que forma a experiência
tem; \arq{salas/} diz se ela está aberta naquela turma; \arq{recursos/} diz se
ela EXISTE naquela rede --- e é pré-condição dos outros dois. Ele expõe a
leitura para qualquer autenticado e a ESCRITA só ao \cd{admin} (\adrr{ADR-88}),
em \emph{controllers} separados; a guarda vive em \arq{comum/}, ao lado da de
liberação, porque é de dentro de transações de outros módulos que ela é chamada.
A superfície dele no hub é a área \rota{/unooria/*}, a única restrita a um papel
que o \cd{gestor\_contrato} não tem.

\textbf{A porta da sala é a exceção que confirma a regra.} \arq{salas/} é módulo de
PLATAFORMA: ele não conhece missão nenhuma pelo nome --- sabe turma, série e
catálogo. E a guarda que as experiências chamam (\cd{exigirMissaoLiberada}) NÃO
mora nele, mora em \arq{comum/liberacao.ts}. Fosse um \emph{provider} do
\arq{salas/}, cada módulo de experiência teria de importá-lo, e uma dependência
esquecida viraria uma porta aberta em silêncio (\adrr{ADR-69}).

\textbf{\arq{equipes/} é o segundo módulo de plataforma no mesmo molde}
(\adrr{ADR-111}). Ele diz QUEM DIVIDE UM COMPUTADOR em cada missão de cada
turma, e a única coisa que sabe de missão é a lista
\cd{EXPERIENCIAS\_COM\_EQUIPE}. O que cada missão FAZ com a equipe --- o painel
contando por computador --- é assunto do módulo dela: \arq{unootube/} e
\arq{unoogpt/} leem \arq{equipes/consulta.ts} (funções livres que recebem o
\cd{tx}, como \arq{comum/liberacao.ts}) e aplicam \cd{colapsarEmUnidades}, de
\arq{equipes/regras.ts}, sem injetar provedor nenhum. O sentido inverso --- a
missão dizendo ao sorteio o lado de cada criança --- é um REGISTRO: o módulo da
missão chama \cd{EquipesService.registrarFonteDeGrupo} na subida, e
\arq{equipes/} nunca importa nada dele. Quem registra hoje é o Unoovista
(\arq{perfis/fonte-de-lado.ts}, H13.4.1), e a fonte recebe a ATIVIDADE, porque o
lado do feed é por atividade: o perfil importado, a afinidade da pergunta e a
redistribuição do professor. No hub, a seção mora na ficha da missão
(\arq{components/equipes/}), e as chamadas passam pelo BFF
(\rota{/api/plataforma/equipes/*}), do lado do cliente porque o professor monta
as equipes na frente da sala.

\textbf{\arq{evidencias/} é o terceiro, e o primeiro que as missões CHAMAM}
(\adrr{ADR-112}). Ele guarda que habilidades da BNCC cada estudante demonstrou
numa atividade de uma missão, rodada a rodada, e não conhece missão nenhuma pelo
nome: o que cada uma sabe de si entra por três pontos. A \textbf{ficha} é dado
estático do módulo (\arq{evidencias/\B fichas/\B <expId>.ts}, no registro
\cd{FICHAS}, que tem uma só: a da exp05ef1b, \adrr{ADR-118}), porque o gabarito
não pode ir para o \emph{bundle}. O \textbf{contexto} --- em que turma, atividade e rodada o
estudante está agora --- é um REGISTRO, como a fonte de lado das equipes: o módulo
da missão chama \cd{EvidenciasService.\B{}registrarContexto(expId, resolver)} na
subida. E a \textbf{marca} é uma chamada, \cd{EvidenciasService.marcar(tx,
\ldots)}, feita DA transação da missão, sem abrir \cd{comTenancy} --- a marca cai
junto se a jogada cair. Por isso o módulo exporta o \emph{service}, que é a
porta das missões, e a dependência vai num sentido só: a missão conhece a
evidência, e \arq{evidencias/} nunca importa módulo de missão. Com \arq{equipes/} o arranjo é o inverso: \arq{evidencias/} importa
o \cd{EquipesModule}, lê \arq{equipes/consulta.ts} (\cd{mapaDeEquipes},
\cd{computadoresDaMontagem}) e registra o recálculo das cópias no gancho
\cd{EquipesService.\B{}aoMudarEquipes}, que o \cd{salvar} e o \cd{desfazer} chamam
dentro da própria transação; \arq{equipes/} não sabe que a evidência existe. As
regras moram em \arq{evidencias/regras.ts}, puras, e os códigos da BNCC vêm de
\arq{curriculo/} (\cd{habilidadePorCodigo}). No hub, o painel
\enquote{Habilidades da BNCC} é genérico, desenhado só quando a resposta diz
\cd{temFicha: true}: entra na página da atividade do Instanooria (que ainda não
tem ficha) e nos resultados da exp05ef1b, onde a atividade é o pacote aprovado
--- é lá que ele aparece hoje. O perfil do estudante é uma página nova, aberta
pela aba Estudantes da turma.

\textbf{Módulo de domínio não conhece módulo de domínio.} \arq{vila/} não importa
\arq{unootube/}. Quando duas experiências precisam se comunicar --- o perfil da
Vila alimentando a largada do Instanooria --- a comunicação é por \emph{tabela de
contrato} explícita (\cd{exp08ef1b\_app\_70\_perfis\_importados}, que carrega os
dois números no nome) e por um \emph{endpoint} de fronteira
(\rota{POST /atividades/:id/perfis/da-vila}). Chamar a tabela por um número só
mentiria sobre quem depende dela.
\end{regra}

\subsection{Componentes do hub}

\begin{figure}[H]
\centering
\fitwidth{%
\begin{tikzpicture}[x=1cm, y=1cm]

  \draw[c4zone] (0.5,-0.5) rectangle (15.5,-10.3);
  \node[c4boundarylbl, fill=unooSurface] at (0.9,-0.35) {Contêiner: Hub (Next.js)};

  \node[camadalbl, anchor=west] at (0.8,-1.2) {Borda};
  \node[c4component, text width=29mm] (mw) at (2.7,-2.3)
    {\texttt{middleware.ts}\desc{sessão + área do papel}};
  \node[c4component, text width=29mm] (papeis) at (6.4,-2.3)
    {\texttt{lib/papeis.ts}\desc{\emph{edge-safe}: papel \seta{} área}};
  \node[c4component, text width=29mm] (authlib) at (10.1,-2.3)
    {\texttt{lib/auth.ts}\desc{NextAuth, sessão}};
  \node[c4component, text width=29mm] (jwks) at (13.6,-2.3)
    {\texttt{/.well-known/jwks.json}\desc{chave pública}};

  \node[camadalbl, anchor=west] at (0.8,-3.9) {Montagem do \emph{embed}};
  \node[c4component, text width=32mm] (embed) at (3.0,-5.1)
    {\texttt{lib/embed-experiencia.ts}\desc{token + manifesto + origem, uma só vez}};
  \node[c4component, text width=28mm] (chave) at (7.2,-5.1)
    {\texttt{lib/chave-token.ts}\desc{assina RS256, TTL 15 min}};
  \node[c4component, text width=28mm] (devsrv) at (11.0,-5.1)
    {\texttt{lib/experiencias-dev.ts}\desc{porta do \emph{dev server}}};
  \node[c4component, text width=24mm] (apiserv) at (14.4,-5.1)
    {\texttt{lib/api-servidor.ts}\desc{BFF}};

  \node[camadalbl, anchor=west] at (0.8,-6.7) {Casca visual e \emph{host} do \texttt{iframe}};
  \node[c4component, text width=30mm] (casca) at (2.9,-7.9)
    {\texttt{CascaPlataforma}\desc{CSS do DS + fontes}};
  \node[c4component, text width=30mm] (header) at (6.7,-7.9)
    {\texttt{HeaderPlataforma}\desc{um só para todos os papéis}};
  \node[c4component, text width=28mm] (disp) at (10.4,-7.9)
    {\texttt{DispositivoFrame}\desc{casca de aparelho}};
  \node[c4component, text width=28mm] (frame) at (13.9,-7.9)
    {\texttt{ExpFrame}\desc{\emph{sandbox} + protocolo v1}};

  \node[c4component, text width=30mm] (rotas) at (4.4,-9.5)
    {Rotas: \texttt{/} (\emph{landing}), \texttt{/privacidade}, \texttt{/termos},
     \texttt{/estudante}, \texttt{/professor/*},
     \texttt{/admin/*}, \texttt{/analise/*}, \texttt{/unooria/*},
     \texttt{/exp/:expId}, \texttt{/dev-embed}};
  \node[c4component, text width=30mm] (bff) at (10.4,-9.5)
    {\texttt{/api/plataforma/[...path]}\desc{\emph{proxy} autenticado para a API}};

  \draw[c4rel] (mw.east) -- (papeis.west);
  \draw[c4rel] (embed.east) -- (chave.west);
  \draw[c4rel] (embed.south) -- (4.4,-8.9);
  \draw[c4rel] (disp.south east) -- (frame.north west);

\end{tikzpicture}}
\caption[C4 nível 3 --- componentes do hub]{C4 nível 3: componentes do hub. A
  montagem do \emph{embed} é \emph{uma só função} usada pelas quatro superfícies
  que embedam (\rota{/exp/[expId]}, \rota{/dev-embed},
  \rota{/professor/plano-de-aula} e \rota{/professor/bem-estar}) --- duplicá-la
  era o caminho mais curto para uma delas passar a assinar o \emph{token} com
  \emph{claims} diferentes.}
\label{fig:c4l3hub}
\end{figure}

\begin{decisao}
\textbf{Um único \emph{header} para todos os papéis, e uma única casca visual.}
Antes, cada área copiava o próprio \emph{header} e a experiência não tinha
nenhum. Hoje \cd{HeaderPlataforma} serve \rota{/estudante}, \rota{/professor},
\rota{/admin}, \rota{/analise} \emph{e} \rota{/exp/*}, e a casca visual saiu do
\emph{layout} do grupo de rotas para \cd{CascaPlataforma}, porque \rota{/exp/*} e
\rota{/dev-embed} estão fora do grupo e também precisam dela. A regra de ouro
passou a valer um nível abaixo: o \emph{header} é \emph{chrome} do hub, desenhado
\emph{acima} do aparelho --- dentro do \cd{iframe} nada da marca institucional
entra.
\end{decisao}

\begin{decisao}
\textbf{O \emph{layout} de cada área é onde moram as regras que valem para a área
inteira.} São três hoje, e nenhuma pertence à página: a \textbf{trilha}
(\adrr{ADR-77}), a \textbf{largura da coluna} (\adrr{ADR-83}) e a
\textbf{fronteira de remontagem por escola ativa} (\adrr{ADR-85}, o
\cd{EscopoDaEscola} que envolve o miolo). As três chegaram ali pelo mesmo
caminho: eram responsabilidade de cada página, e responsabilidade de página se
esquece --- 15 telas sem \emph{breadcrumb}, quatro larguras de coluna sem razão
registrada e quatro telas que não reagiam à troca de escola. \textbf{Área nova
precisa de \emph{layout} novo}, e é o \emph{layout} que faz as três coisas; sem
ele a área nasce com os três defeitos, como \rota{/admin} nasceu.
\end{decisao}

\begin{decisao}
\textbf{O chat do agente de IA do professor é do \emph{layout} de
\rota{/professor}, não de uma página} (\adrr{ADR-126}). O botão flutuante
(\cd{ChatLauncher} do DS) e o painel não modal (\cd{ChatPanel}, compacto ou
expandido) moram em \arq{components/agente-professor/}, montados pelo
\emph{layout} da área só quando o produto \cd{agente-professor} está ligado
para a escola --- o produto não tem \emph{card} na \emph{home}: a superfície é o
botão. Duas peças fogem do desenho do resto do hub. A \textbf{rota de
\emph{stream}} (\rota{/api/plataforma/agente-professor/\ldots}) é própria, fora do
BFF genérico, porque ele tem teto de 180\,s e devolve o corpo sem os cabeçalhos
de \emph{stream} (\adrr{ADR-128}); a sessão e o \emph{token} curto são os mesmos
(\arq{lib/bff.ts}). Listar, criar, ler e apagar seguem pelo BFF. E o \textbf{\emph{markdown} vira HTML no navegador}, com
\cd{react-markdown} sem HTML cru, porque o texto chega em pedaços
(\adrr{ADR-129}). O painel leva \cd{data-clarity-mask}.
\end{decisao}

\begin{decisao}
\textbf{A Ajuda com avatar é uma janela do menu da conta, não uma rota}
(\adrr{ADR-132}). O item \enquote{Ajuda} é o segundo do menu do
\cd{HeaderPlataforma}, e o \emph{header} só o desenha para quem não é estudante e
com o produto \cd{ajuda} ligado --- o mesmo \emph{header} serve \rota{/exp/*},
então esconder o item é dele; quem tranca é a API. O clique abre um \cd{Modal}
do DS (\arq{components/\B ajuda/\B modal-da-ajuda.tsx}): o atendente em vídeo em
cima e a conversa escrita embaixo, com \cd{ChatMessage}, \cd{ChatComposer},
\cd{TypingIndicator}, \cd{ScrollBar} e \cd{AlertBox}. O \cd{<video>} é o único
elemento cru, e é mídia, não controle: o DS não publica \emph{player} ao vivo e o
SDK do \emph{vendor} pede um \cd{HTMLMediaElement}. O ciclo da sala mora no
\emph{hook} \arq{components/\B ajuda/\B use-sessao-da-ajuda.ts}, que carrega o SDK
do \emph{vendor} só por \cd{import()} (cerca de 1\,MB de LiveKit que nenhuma outra
tela usa), começa com o microfone desligado e guarda o \emph{token} da sala só na
\emph{closure} (regra 4). A sala abre na \textbf{montagem} do conteúdo do modal e
fecha na \textbf{desmontagem}: renderizar o conteúdo com o modal fechado abriria
sessões, e cobraria minutos, sem ninguém ver. \arq{lib/\B ajuda/} tem as duas
peças puras, com teste: \arq{transporte.ts} (abrir e encerrar pelo BFF genérico,
\rota{/api/plataforma/ajuda/\ldots}) e \arq{conversa.ts}, o redutor que monta a
conversa das \emph{transcrições} que o \emph{vendor} devolve pela sala --- o texto
não passa pelo hub. Nada é guardado, e o modal leva \cd{data-clarity-mask}.
\end{decisao}

\begin{nota}
\textbf{A raiz é a \emph{landing} pública, e continua roteando quem tem sessão}
(\adrr{ADR-40}). \rota{/} sem sessão renderiza \cd{components/landing/} ---
porte da tela desenhada no Stitch, com a estrutura, a copy e o movimento do
modelo (navegação fixa em vidro, \emph{shader} WebGL de fundo, revelação ao
rolar) e a cor, a tipografia, os ícones, o logo e os botões do DS. Com sessão,
\rota{/} redireciona para a \emph{home} do papel (ou para \rota{/trocar-senha}):
o \cd{login-form} e a troca de senha fazem \cd{router.replace("/")} e dependem
desse roteamento para cair na área certa sem o \emph{client} conhecer o papel. A
\emph{landing} está fora do \cd{matcher} do \emph{middleware} e não lê nada da
sessão. Os dois pedaços com ciclo de vida (\cd{FundoAnimado},
\cd{RevelarAoRolar}) são \emph{client components}; o resto é servidor.
\end{nota}

\begin{nota}
\textbf{As páginas legais são a mesma superfície pública da \emph{landing}}
(\adrr{ADR-70}). \rota{/privacidade} e \rota{/termos} completam o conjunto
público do hub --- com \rota{/}, \rota{/login} e o JWKS, é tudo que existe sem
sessão. Elas são alcançadas pelo rodapé da \emph{landing}, que até então
oferecia os dois links com \cd{href="\#"}, e por isso reaproveitam a casca dela:
\cd{components/legal/} importa \cd{landing.css} e monta o mesmo
\cd{FundoAnimado}, a mesma navegação em vidro e a mesma escala tipográfica.
Duplicar aquilo criaria a segunda cópia, que diverge na primeira mudança de
marca.

O conteúdo entra como \emph{dado} (uma lista de seções com \cd{id},
\cd{titulo} e \cd{conteudo}), não como marcação livre: é o que mantém o sumário
lateral em sincronia com o texto --- numeração, âncora e ordem saem da mesma
lista que renderiza as seções, então não existe o estado em que o índice promete
um item que o corpo não tem. O único \emph{client component} é o
\cd{SumarioAtivo}, que acende o item em leitura.
\end{nota}

\section{C4 nível 4 --- Código}
\label{sec:c4l4}

O nível 4 do C4 raramente vale a pena desenhar por inteiro. Vale aqui para
\emph{três} lugares, porque neles o desenho de tipos \emph{é} a arquitetura.

\subsection{O contrato da ponte e o \emph{guard} de \emph{tenancy}}

\begin{figure}[H]
\centering
\fitwidth{%
\begin{tikzpicture}[x=1cm, y=1cm, node distance=6mm]

  \node[classe, text width=48mm] (bridge) at (0,0)
    {\etitulo{connectToHub}
     \nodepart{two}
     \texttt{hubOrigin: string}\\
     \texttt{timeoutMs = 8000}\\
     \texttt{READY\_RETRY\_MS = 400}
     \nodepart{three}
     valida \texttt{ev.origin} \seta{} descarta\\
     reenvia \texttt{exp:ready} até responder\\
     compara \texttt{protocolVersion} maior\\
     valida pós-\emph{init} com Zod};

  \node[classe, text width=44mm, right=14mm of bridge] (types)
    {\etitulo{@edu/sdk/types}
     \nodepart{two}
     \texttt{PROTOCOL\_VERSION = 2}\\
     \texttt{HubInitSchema}\\
     \texttt{HubParaExpSchema}\\
     \texttt{ExpParaHubSchema}
     \nodepart{three}
     uniões discriminadas por \texttt{type}\\
     mensagem que não casa é\\
     descartada com aviso};

  \node[classe, text width=46mm, below=13mm of bridge] (prismas)
    {\etitulo{PrismaService}
     \nodepart{two}
     \texttt{extends PrismaClient}
     \nodepart{three}
     \texttt{comTenancy(fn, escolaId?, opts?)}\\
     \quad abre \texttt{\$transaction}\\
     \quad \texttt{set\_config('app.escola\_id', ..., true)}\\
     \quad executa \texttt{fn(tx)}\\
     sem escola \seta{} \textbf{lança}};

  \node[classe, text width=42mm, right=14mm of prismas] (ctx)
    {\etitulo{tenancy-context}
     \nodepart{two}
     \texttt{AsyncLocalStorage<string>}
     \nodepart{three}
     \texttt{comEscolaNoContexto(id, fn)}\\
     \texttt{escolaDoContexto(): string?}};

  \node[classe, text width=44mm, below=13mm of prismas] (leitura)
    {\etitulo{LinhaDeLeitura (união)}
     \nodepart{two}
     \texttt{\{ grau: 'scored', numero: number \}}\\
     \texttt{\{ grau: 'indicative', numero: null \}}\\
     \texttt{\{ grau: 'zero', numero: null \}}\\
     \texttt{\{ grau: 'none', numero: null \}}
     \nodepart{three}
     o hub \emph{não pode} imprimir número\\
     onde não há medida --- o tipo\\
     não o oferece};

  \node[actnote, text width=52mm, right=14mm of leitura] (nota)
    {\textbf{Por que nível 4 aqui:} nos três casos a regra é imposta pelo
     \emph{tipo}, não por revisão de código. Comentário avisa; tipo impede.
     \texttt{comTenancy} sem escola \emph{lança} em vez de consultar sem filtro;
     mensagem inválida é descartada em vez de despachada crua; leitura sem
     medida não tem campo numérico para ser lido.};

  \draw[depend] (bridge) -- (types) node[cxlabel, midway, above]{usa};
  \draw[depend] (prismas) -- (ctx) node[cxlabel, midway, above]{lê};
  \draw[actnoteline] (leitura) -- (nota);

\end{tikzpicture}}
\caption[C4 nível 4 --- tipos que carregam a regra]{C4 nível 4: os três pontos em
  que a regra arquitetural está no sistema de tipos. \cd{comTenancy} é a única
  porta do banco e falha alto sem escola; o contrato da ponte descarta o que não
  valida; a leitura de valores não oferece número onde não há medida.}
\label{fig:c4l4}
\end{figure}

\subsection{O gabarito de um módulo de domínio}

Todo módulo de experiência segue a mesma forma. Não é convenção estética: é o que
permite testar a regra sem banco e trocar a persistência sem tocar a regra.

\begin{lstlisting}[language=tsunoo, caption={Forma canônica de um módulo de domínio de experiência}, label={lst:modulo}]
apps/api/src/<slug>/
  catalogo.ts     // DADOS: itens, textos, gabarito. Constantes puras.
  regras.ts       // REGRA: funcoes puras. Sem Prisma, sem fetch, sem Date.
  <slug>.service.ts   // PERSISTENCIA: le com lock, aplica a regra, grava.
                      // Toda query dentro de comTenancy e filtrando escola_id.
  <slug>.controller.ts        // casca HTTP do papel "estudante"
  <slug>-professor.controller.ts  // casca HTTP do papel "professor"
  <slug>.module.ts

// A separacao de controllers POR PAPEL e deliberada: um @Roles misto com
// excecao por metodo e o caminho mais curto para um GET de professor herdar
// silenciosamente a permissao de um POST de estudante.
\end{lstlisting}

\begin{nota}
A consequência prática dessa forma aparece nos testes: \arq{vila-regras.spec.ts},
\arq{unootube-regras.spec.ts} e \arq{bemestar-regras.spec.ts} rodam
\emph{milissegundos} sem PostgreSQL, porque exercitam \arq{regras.ts} puro. Os
testes que precisam de banco (\arq{*-postgres.spec.ts}, \arq{tenancy.spec.ts})
são separados e cobrem exatamente o que não dá para afirmar sem ele:
\emph{tenancy}, \emph{lock} e concorrência.
\end{nota}

\section{Estrutura do \emph{monorepo}}
\label{sec:monorepo}

\begin{figure}[H]
\centering
\fitwidth{%
\begin{tikzpicture}[x=1cm, y=1cm]

  \node[c4component, text width=28mm, fill=unooBlueWash] (sdk) at (7.5,0)
    {\texttt{packages/sdk}\desc{Zod + \emph{fetch}. Zero dependência pesada}};
  \node[c4component, text width=28mm, fill=unooBlueWash] (ui) at (12.0,0)
    {\texttt{packages/ui}\desc{ponte para o \emph{design system}}};

  \node[c4container, text width=26mm] (hub) at (2.0,-3.2) {\texttt{apps/hub}};
  \node[c4container, text width=26mm] (api) at (5.6,-3.2) {\texttt{apps/api}};
  \node[bundle, text width=24mm] (e01) at (9.0,-3.2) {\texttt{exp08ef1b}};
  \node[bundle, text width=24mm] (e02) at (11.9,-3.2) {\texttt{exp04ef1b}};
  \node[bundle, text width=24mm] (e70) at (14.8,-3.2) {\texttt{app-70}};
  \node[bundle, text width=24mm] (e90) at (9.0,-5.2) {\texttt{app-90}};
  \node[bundle, text width=24mm] (e91) at (11.9,-5.2) {\texttt{app-91}};
  \node[bundle, text width=24mm, fill=unooMuted, draw=unooInkSoft] (tpl) at (14.8,-5.2)
    {\texttt{\_template}};

  \node[c4systemext, text width=34mm] (ds) at (12.0,2.6)
    {\texttt{@unooria/design-system}\desc{npm público, repo \texttt{storybook-unooria}}};

  \draw[c4rel] (ds.south) -- (ui.north) node[cxlabel, midway, right=2pt]{versão publicada};
  \draw[c4rel] (sdk.south) -- (hub.north east);
  \draw[c4rel] (sdk.south) -- (e01.north);
  \draw[c4rel] (sdk.south) -- (e02.north);
  \draw[c4rel] (sdk.south) -- (e70.north);
  \draw[c4rel] (sdk.south) -- (9.0,-4.4) -- (e90.north);
  \draw[c4rel] (sdk.south) -- (11.9,-4.4) -- (e91.north);
  \draw[c4rel] (ui.south) -- (hub.north east);
  \draw[c4rel] (ui.south) -- (e01.north east);
  \draw[c4rel] (ui.south) -- (e02.north east);
  \draw[c4rel] (ui.south) -- (e90.north east);
  \draw[c4rel] (ui.south) -- (e91.north east);

\end{tikzpicture}}
\caption[Grafo de dependências do monorepo]{Grafo de dependências do
  \emph{monorepo}, com uma amostra das experiências. A SDK é consumida pelo hub e
  por toda experiência --- é por isso que uma mudança nela precisa testar todas
  no mesmo \emph{pull request}, e por isso o \emph{monorepo} (\adrr{ADR-02}). A API
  não importa pacote nenhum do \emph{workspace}: valida o que recebe com os
  próprios \emph{schemas} (Figura~\ref{fig:uml-pacotes}). A app-70 não depende de \arq{packages/ui}: é
  jogo em \emph{canvas} 3D, sem DOM de interface, e é a exceção registrada em ADR
  à regra do \emph{design system}.}
\label{fig:monorepo}
\end{figure}

\begin{tabularx}{\linewidth}{@{}L{34mm}X@{}}
\toprule
\headrow \thd{Diretório} & \thd{Conteúdo}\\
\midrule
\arq{apps/hub} & Aplicação Next.js: rotas por papel, \emph{embed}, BFF, JWKS.\\
\zebra \arq{apps/api} & Aplicação NestJS: 34 módulos, Prisma, 40 migrações
  (contagem de 2026-09-25).\\
\arq{apps/experiencias/*} & Uma pasta por experiência, mais \arq{\_template}.\\
\zebra \arq{packages/sdk} & Contrato da ponte e telemetria.\\
\arq{packages/ui} & Ponte do \emph{design system} e CSS da plataforma.\\
\zebra \arq{docs/} & Contratos com infraestrutura e registro de decisões.\\
\arq{scripts/} & \emph{Scaffolding}, chave de assinatura, importação de
  referências, a suíte de ponta a ponta (\arq{scripts/e2e/}), \emph{probe} de
  \emph{pointer lock}.\\
\zebra \arq{wso2/} & \emph{Compose} próprio do provedor de identidade, com ciclo
  de vida separado.\\
\arq{documentacao-latex/} & Este documento.\\
\bottomrule
\end{tabularx}

A orquestração é Turborepo~\cite{turborepo2026}, com \emph{pipeline} declarado
em \arq{turbo.json}: \cd{build} depende do \cd{build} das dependências;
\cd{typecheck} e \cd{test} também. É isso que permite ao CI rodar apenas o que
mudou:

\begin{lstlisting}[language=bashunoo, caption={O comando que o CI executa}]
pnpm turbo typecheck lint test build --filter='...[HEAD^]'
\end{lstlisting}

\section{Fluxos principais}
\label{sec:fluxos}

\subsection{Curadoria de um bloco de \emph{feed} (exp08ef1b)}

Este é o fluxo com mais concorrência da plataforma e o que mais exige do banco.

\begin{figure}[H]
\centering
\fitwidth{%
\begin{tikzpicture}[x=1cm, y=1cm]

  \raia{1.4}{-10.4}{Experiência}{exp08ef1b no \texttt{iframe}}
  \raia{5.4}{-10.4}{\texttt{runs.controller}}{NestJS}
  \raia{9.4}{-10.4}{\texttt{motor/servico}}{regra + persistência}
  \raiaext{13.4}{-10.4}{PostgreSQL}{Cloud SQL}

  \msgs{1.4}{5.4}{-1.1}{\texttt{POST /runs/:id/eventos} (\emph{dwell}, curtidas)}
  \msgs{5.4}{9.4}{-1.9}{\texttt{aplicarEventos}}
  \msgs{9.4}{13.4}{-2.7}{\texttt{comTenancy}: \texttt{set\_config('app.escola\_id')}}
  \msgs{9.4}{13.4}{-3.5}{\texttt{pg\_advisory\_xact\_lock(runId)}}
  \msgr{13.4}{9.4}{-4.3}{estado da \texttt{run}}

  \msgself{9.4}{-5.1}{\texttt{curarBloco}: puro e determinístico}
  \msgs{9.4}{13.4}{-6.3}{\texttt{UPDATE engine\_state} + posts entregues}
  \msgr{13.4}{9.4}{-7.1}{\texttt{COMMIT} (libera o \emph{lock})}
  \msgr{9.4}{5.4}{-7.9}{bloco + módulos + diagnóstico}
  \msgself{5.4}{-8.7}{\texttt{paraFeedPost}: remove metadados de curadoria}
  \msgr{5.4}{1.4}{-9.9}{\texttt{GET /runs/:id/feed} --- \textbf{sem} lado, perfil nem intensidade}

  \node[actnote, text width=44mm, anchor=north west] at (15.0,-1.0)
    {\textbf{Por que o \emph{lock}:} dois pedidos simultâneos de bloco (duas abas,
     um \emph{reload} no meio do \emph{fetch}, um \emph{retry}) se sobrescreviam
     e um bloco inteiro desaparecia do estado --- posts já entregues voltavam a
     ser servidos e a série por bloco do dossiê ficava corrompida. Transação com
     \texttt{timeout} de \SI{20}{\second}.};
  \node[actnote, text width=44mm, anchor=north west] at (15.0,-6.0)
    {\textbf{Por que o sanitizador:} \texttt{lado}, \texttt{perfil} e \texttt{intensidade}
     são \emph{a variável independente do experimento}. Se chegassem ao cliente,
     a experiência poderia exibi-los --- e o estudante saberia de que lado o
     algoritmo o colocou antes da revelação.};

\end{tikzpicture}}
\caption[Sequência: curadoria de um bloco de feed]{Sequência da curadoria de um
  bloco de \emph{feed}. A regra (\cd{curarBloco}) é pura; o \emph{service} lê com
  \emph{lock}, aplica e grava; o \emph{controller} sanitiza. As três
  responsabilidades ficam em arquivos diferentes de propósito.}
\label{fig:seq-feed}
\end{figure}

\subsection{Estados de uma geração assíncrona}

Duas rotinas da plataforma são assíncronas e longas: a geração de \emph{feed} da
exp08ef1b e a geração de plano de aula da app-90. Ambas seguem a mesma máquina de
estados, com \emph{claim} atômico em vez de fila externa.

\begin{figure}[H]
\centering
\fitwidth{%
\begin{tikzpicture}[x=1cm, y=1cm]

  \node[actstart] (ini) at (0,0) {};
  \node[estado, right=10mm of ini] (pend) {\texttt{pendente}};
  \node[estado, right=16mm of pend] (ger) {\texttt{gerando}};
  \node[estadofim, right=16mm of ger] (ok) {\texttt{pronto}};
  \node[estadoerr, below=13mm of ger] (err) {\texttt{erro}};
  \node[actend, right=12mm of ok] (fim) {};

  \draw[trans] (ini) -- (pend) node[translbl, midway, above]{\texttt{POST} (202)};
  \draw[trans] (pend) -- (ger)
    node[translbl, midway, above, align=center]{\emph{claim} atômico\\\texttt{updateMany} por \texttt{status}};
  \draw[trans] (ger) -- (ok) node[translbl, midway, above]{gravou o artefato};
  \draw[trans] (ger) -- (err)
    node[translbl, midway, right=2pt, align=left]{falha do \emph{provider}\\ou \texttt{finishReason=MAX\_TOKENS}};
  \draw[trans] (ok) -- (fim);
  \draw[trans] (ger.north) to[out=60, in=120, looseness=1.6] (pend.north);
  \node[translbl] at (7.0,1.35) {órfã ($>$ 12 min sem \emph{worker}) --- até 2 retomadas};

  \node[actnote, text width=54mm, anchor=north west] at (0,-3.0)
    {O \emph{claim} atômico (\texttt{UPDATE ... WHERE status = 'pendente'} com
     contagem de linhas afetadas) substitui uma fila externa: duas instâncias do
     Cloud Run podem tentar, e só uma ganha a linha. É a menor peça possível que
     resolve o problema --- e não exige operar um serviço a mais
     (\adrr{ADR-17}).};

\end{tikzpicture}}
\caption[Estados de uma geração assíncrona]{Máquina de estados das gerações
  assíncronas (exp08ef1b e app-90). O \emph{claim} atômico dispensa fila externa; a
  retomada de órfã é disparada pelo \emph{polling} do próprio painel, não por um
  \emph{scheduler}.}
\label{fig:estados-geracao}
\end{figure}

\begin{limitacao}
Não há retomada no \emph{boot}. Se todas as instâncias morrerem e ninguém abrir o
painel, uma geração órfã fica parada até que alguém abra. É aceito porque o
\emph{polling} do painel é justamente o momento em que alguém se importa com
aquela geração --- e porque a alternativa é um \emph{scheduler}, que é uma peça
nova para operar. O gatilho de reavaliação é a primeira reclamação real de
geração parada.
\end{limitacao}
